\documentclass{beamer}
\usetheme{Madrid}
\usecolortheme{default}

\title{Cross-Field --- Representative Vector}
\subtitle{Domain-Partition Algorithm for Turbine Surfaces}
\author{Domain-Partition 3D Project}
\date{\today}
\institute{Block-Structured Meshing Pipeline}

\begin{document}

\begin{frame}
  \titlepage
\end{frame}

\begin{frame}{Cross-Field $\rightarrow$ Representative Vector}
  \begin{columns}
    \begin{column}{0.55\textwidth}
      \begin{itemize}
        \item Each cross has 4 symmetric representatives (rotated by $\pi/2$)
        \item Pick the one closest to the desired direction
        \item Enables streamline integration along smooth field lines
      \end{itemize}
    \end{column}
    \begin{column}{0.45\textwidth}
      \centering
      \footnotesize
      \textbf{Vector mapping:}
      \vspace{0.3cm}
      \[
      \mathbf{v}(x) = \arg\max_{\mathbf{c} \in \text{cross}(x)} \langle \mathbf{c}, \mathbf{d}(x) \rangle
      \]
      \vspace{0.2cm}
      where $\mathbf{d}(x)$ is the desired direction.
    \end{column}
  \end{columns}

  \vspace{0.5cm}
  \centering
  \footnotesize
  Citation: Kn\"{o}ppel et al., "Globally optimal direction fields," SIGGRAPH 2013.
\end{frame}

\note{
  Speaker notes:
  This is slide 4 of an 11-slide presentation on the domain-partition algorithm.
  A 4-RoSy cross-field assigns four symmetric directions at every point on the surface.
  To integrate streamlines, we must pick one representative vector from the cross.
  The mapping chooses the representative that maximizes the inner product with the desired direction,
  ensuring smooth transitions and avoiding abrupt jumps between adjacent crosses.
  This step is essential before streamline integration and singularity detection can proceed.
}

\end{document}
